% GENERATED FILE. Edit canonical lessons, then run npm run generate:books.
% canonical-source-hash: fnv1a64:1ac95cdce79fad54
% canonical-lessons: FA-C84-hame, FA-C84-kasi, FA-C84-chizi, FA-C84-hich, FA-C84-khod

\chapter{All, Someone, Something, None, Oneself}
\label{ch:fa-all-someone-something-none-oneself}
\hlchaptermodality{84}

\begin{chapteropening}
\textbf{By the end of this chapter:} \emph{I can say all, someone, something, none and oneself.}

Complete the last lesson of chapter 84: I can say all, someone, something, none and oneself.
\end{chapteropening}

\section[\texorpdfstring{\fa{همه} (hame)}{همه (hame)}]{\fa{همه} (hame) — all, everyone}

\label{lesson:FA-C84-hame}



\begin{quote}
Before the new one: say the Persian for such, then the Persian for every.
\end{quote}

\subsection*{You'll want to know: \fa{همه}}
\textbf{\fa{همه}} — \emph{hame} — \textquotedblleft{}all, everyone\textquotedblright{}.

\begin{cousinweb}
\emph{hame} is a native Persian word, inherited rather than borrowed: it continues Middle Persian \emph{hamāg}, \textquotedblleft{}all\textquotedblright{} (EtymDB).
\end{cousinweb}

\begin{grammarlens}[title={What you've built}]
One more word for this chapter.
\end{grammarlens}

\subsection*{Your turn}
\begin{itemize}
\raggedright
  \item \emph{Say it:} \emph{hame}
  \item \emph{Say it:} \emph{hame}, once more
  \item \emph{Say it:} say \emph{hame}
  \item \emph{Recall:} say the Persian for such, then the Persian for every, then say \emph{hame} again
\end{itemize}

\begin{tcolorbox}[breakable,colback=teal!4,colframe=teal!35!black,title={Before you move on}]
What does \fa{همه} mean? (\textquotedblleft{}All, everyone\textquotedblright{}.) Say it once more.
\end{tcolorbox}
\section[\texorpdfstring{\fa{کسی} (kasi)}{کسی (kasi)}]{\fa{کسی} (kasi) — someone}

\label{lesson:FA-C84-kasi}



\begin{quote}
Before the new one: say the Persian for every, then the Persian for all.
\end{quote}

\subsection*{You'll want to know: \fa{کسی}}
\textbf{\fa{کسی}} — \emph{kasi} — \textquotedblleft{}someone\textquotedblright{}.

\begin{cousinweb}
\emph{kasi} is \emph{kas}, \textquotedblleft{}person\textquotedblright{}, with the ending \emph{-i}, \textquotedblleft{}a, some\textquotedblright{}: \textquotedblleft{}some person\textquotedblright{}. \emph{kas} is native, usually traced to Middle Persian \emph{kas}.
\end{cousinweb}

\begin{grammarlens}[title={What you've built}]
One more word for this chapter.
\end{grammarlens}

\subsection*{Your turn}
\begin{itemize}
\raggedright
  \item \emph{Say it:} \emph{kasi}
  \item \emph{Say it:} \emph{kasi}, once more
  \item \emph{Say it:} say \emph{kasi}
  \item \emph{Recall:} say the Persian for every, then the Persian for all, then say \emph{kasi} again
\end{itemize}

\begin{tcolorbox}[breakable,colback=teal!4,colframe=teal!35!black,title={Before you move on}]
What does \fa{کسی} mean? (\textquotedblleft{}Someone\textquotedblright{}.) Say it once more.
\end{tcolorbox}
\section[\texorpdfstring{\fa{چیزی} (chizi)}{چیزی (chizi)}]{\fa{چیزی} (chizi) — something}

\label{lesson:FA-C84-chizi}



\begin{quote}
Before the new one: say the Persian for all, then the Persian for someone.
\end{quote}

\subsection*{You'll want to know: \fa{چیزی}}
\textbf{\fa{چیزی}} — \emph{chizi} — \textquotedblleft{}something\textquotedblright{}.

\begin{cousinweb}
\emph{chizi} is \emph{chiz}, \textquotedblleft{}thing\textquotedblright{}, with the ending \emph{-i}, \textquotedblleft{}a, some\textquotedblright{}: \textquotedblleft{}some thing\textquotedblright{}. \emph{chiz} is native, usually traced to Middle Persian \emph{čiš}.
\end{cousinweb}

\begin{grammarlens}[title={What you've built}]
One more word for this chapter.
\end{grammarlens}

\subsection*{Your turn}
\begin{itemize}
\raggedright
  \item \emph{Say it:} \emph{chizi}
  \item \emph{Say it:} \emph{chizi}, once more
  \item \emph{Say it:} say \emph{chizi}
  \item \emph{Recall:} say the Persian for all, then the Persian for someone, then say \emph{chizi} again
\end{itemize}

\begin{tcolorbox}[breakable,colback=teal!4,colframe=teal!35!black,title={Before you move on}]
What does \fa{چیزی} mean? (\textquotedblleft{}Something\textquotedblright{}.) Say it once more.
\end{tcolorbox}
\section[\texorpdfstring{\fa{هیچ} (hich)}{هیچ (hich)}]{\fa{هیچ} (hich) — none, no}

\label{lesson:FA-C84-hich}



\begin{quote}
Before the new one: say the Persian for someone, then the Persian for something.
\end{quote}

\subsection*{You'll want to know: \fa{هیچ}}
\textbf{\fa{هیچ}} — \emph{hich} — \textquotedblleft{}none, no\textquotedblright{}.

\begin{cousinweb}
\emph{hich} is a native Persian word, inherited rather than borrowed: it continues Middle Persian \emph{ēč}, \textquotedblleft{}any, none\textquotedblright{} (EtymDB).
\end{cousinweb}

\begin{grammarlens}[title={What you've built}]
One more word for this chapter.
\end{grammarlens}

\subsection*{Your turn}
\begin{itemize}
\raggedright
  \item \emph{Say it:} \emph{hich}
  \item \emph{Say it:} \emph{hich}, once more
  \item \emph{Say it:} say \emph{hich}
  \item \emph{Recall:} say the Persian for someone, then the Persian for something, then say \emph{hich} again
\end{itemize}

\begin{tcolorbox}[breakable,colback=teal!4,colframe=teal!35!black,title={Before you move on}]
What does \fa{هیچ} mean? (\textquotedblleft{}None, no\textquotedblright{}.) Say it once more.
\end{tcolorbox}
\section[\texorpdfstring{\fa{خود} (khod)}{خود (khod)}]{\fa{خود} (khod) — oneself}

\label{lesson:FA-C84-khod}



\begin{quote}
Before the new one: say the Persian for something, then the Persian for none.
\end{quote}

\subsection*{You'll want to know: \fa{خود}}
\textbf{\fa{خود}} — \emph{khod} — \textquotedblleft{}oneself\textquotedblright{}.

\begin{cousinweb}
\emph{khod} is a native Persian word, inherited rather than borrowed; it is usually traced to Middle Persian \emph{xwad}, \textquotedblleft{}self\textquotedblright{}. (EtymDB also lists \emph{xōd}, \textquotedblleft{}helmet\textquotedblright{}, a different word spelled the same.)
\end{cousinweb}

\begin{grammarlens}[title={What you've built}]
One more word for this chapter.
\end{grammarlens}

\subsection*{Your turn}
\begin{itemize}
\raggedright
  \item \emph{Say it:} \emph{khod}
  \item \emph{Say it:} \emph{khod}, once more
  \item \emph{Say it:} say \emph{khod}
  \item \emph{Recall:} say the Persian for something, then the Persian for none, then say \emph{khod} again
\end{itemize}

\begin{tcolorbox}[breakable,colback=teal!4,colframe=teal!35!black,title={Before you move on}]
What does \fa{خود} mean? (\textquotedblleft{}Oneself\textquotedblright{}.) Say it once more.
\end{tcolorbox}
